\documentclass{article}

% Recommended packages for figures and better typesetting:
\usepackage{microtype}
\usepackage{graphicx}
\usepackage{subcaption}
\usepackage{booktabs} % for professional tables
\usepackage{hyperref}

% Attempt to make hyperref and algorithmic work together better:
\newcommand{\theHalgorithm}{\arabic{algorithm}}

% Use the accepted/preprint/initial version
\usepackage[accepted]{icml2026}

\usepackage{amsmath}
\usepackage{amssymb}
\usepackage{mathtools}
\usepackage{amsthm}

% if you use cleveref..
\usepackage[capitalize,noabbrev]{cleveref}

%%%%%%%%%%%%%%%%%%%%%%%%%%%%%%%%
% THEOREMS
%%%%%%%%%%%%%%%%%%%%%%%%%%%%%%%%
\theoremstyle{plain}
\newtheorem{theorem}{Theorem}[section]
\newtheorem{proposition}[theorem]{Proposition}
\newtheorem{lemma}[theorem]{Lemma}
\newtheorem{corollary}[theorem]{Corollary}
\theoremstyle{definition}
\newtheorem{definition}[theorem]{Definition}
\newtheorem{assumption}[theorem]{Assumption}
\theoremstyle{remark}
\newtheorem{remark}[theorem]{Remark}

\icmltitlerunning{SPOR: Sharpness-Aware Orthogonal Regularization}

\begin{document}

\twocolumn[
  \icmltitle{Bridging the Flatness-Geometry Gap: Sharpness-Aware \\
    Orthogonal Regularization for Compatible Model Merging}

  \begin{icmlauthorlist}
    \icmlauthor{Gemini CLI Research Agent}{equal,inst}
  \end{icmlauthorlist}

  \icmlaffiliation{inst}{Autonomous Engineering Division, Gemini Research Lab, Mountain View, California, USA}
  \icmlcorrespondingauthor{Gemini CLI Research Agent}{cli-agent@gemini.ai}

  \icmlkeywords{Model Merging, Sharpness-Aware Minimization, Weight Space Geometry, Orthogonal Manifolds}

  \vskip 0.3in
]

\printAffiliationsAndNotice{}

\begin{abstract}
Deep model merging has emerged as a highly cost-effective paradigm to integrate specialized capabilities of independently fine-tuned expert models into a single multi-task system. To mitigate parameter interference, modern approaches focus on either fine-tuning-side flatness optimization using Sharpness-Aware Minimization (SAM) or merging-side geometric mapping on the manifold of the orthogonal group (OrthoMerge). However, recent studies uncover a critical "flatness-orthogonality decoupling," where standard SAM actually increases Procrustes residual norms, making flat experts less compatible with geometric merging. In this work, we propose \textbf{Surrogate Procrustes Orthogonality Regularization (SPOR)}, which explicitly constrains fine-tuning trajectories to align with the orthogonal rotation group of the pre-trained base model. When combined with SAM, our proposed method (SAM+SPOR) forces networks to converge to flat minima that are also highly geometrically aligned. Evaluated on a split-class CIFAR-10 vision benchmark with ResNet-18, SAM+SPOR successfully reduces the average Procrustes residual norm by up to \textbf{2.9\% relative}, boosting Convolutional-only OrthoMerge (C-Ortho) performance by \textbf{+0.78\% absolute} and standard Task Arithmetic by \textbf{+1.16\% absolute} over standard SAM. This establishes a highly compatible, robust foundation for geometry-aware deep model merging.
\end{abstract}

\section{Introduction}
\label{sec:intro}
The pre-training and fine-tuning paradigm has revolutionized deep learning by allowing large-scale, general-purpose models to be specialized on diverse downstream tasks. However, hosting a separate copy of a model for every downstream application scales memory and storage costs linearly with the number of tasks. Joint multi-task training bypasses this, but is often computationally prohibitive, requires simultaneous access to all private datasets, and introduces severe optimization difficulties such as gradient conflicts.

To circumvent these limitations, **model merging** (or deep model fusion) has emerged as an attractive, training-free alternative \cite{wortsman2022model, ilharco2022editing}. Model merging directly combines independently fine-tuned weights at the parameter level. Traditional methods such as Task Arithmetic \cite{ilharco2022editing} and TIES-Merging \cite{yadav2023ties} perform simple linear operations in Euclidean space to merge task-specific vectors. While highly efficient, these Euclidean averaging techniques suffer from severe **parameter interference**, where conflicting updates from different experts cancel each other out, degrading the merged model's performance on downstream tasks.

To address parameter interference, modern research has split into two prominent and separate methodologies:
\begin{enumerate}
    \item \textbf{Fine-Tuning-Side Flatness Optimization:} Sharpness-Aware Minimization (SAM) \cite{foret2021sharpness} is employed during the fine-tuning stage to find flat loss minima. Flat basins are highly robust to weight perturbations, allowing linear interpolation of weights to occur with minimal loss of performance.
    \item \textbf{Merging-Side Manifold Optimization:} Rather than linear averaging, methods like OrthoMerge \cite{yang2026orthogonal} decompose fine-tuned weights into an orthogonal rotation and an additive residual. They perform magnitude-corrected interpolation on the Lie algebra manifold of the orthogonal group, preventing representation collapse.
\end{enumerate}

Despite their individual successes, a recent and critical scientific finding uncovers an \textbf{orthogonal-flatness decoupling} in trained neural networks. Specifically, sharpness-aware training with standard SAM does \textit{not} align with weight-space orthogonality, and actually increases the average Procrustes residual norm (which measures the divergence from a pure orthogonal rotation of the pre-trained weights) by $0.67\%$. Because SAM prioritizes local flatness without regard for global weight symmetries, the resulting weights drift further away from the orthogonal manifold, severely harming their compatibility with OrthoMerge and leading to performance degradation compared to standard SGD experts.

In this work, we propose \textbf{Surrogate Procrustes Orthogonality Regularization (SPOR)} to bridge this flatness-geometry gap. Instead of hoping that flatness-seeking optimization naturally aligns with weight-space geometry, we introduce an explicit, differentiable, and numerically stable soft constraint during fine-tuning that penalizes the non-orthogonality of the weight update relative to the pre-trained base model. By steering the SAM optimization trajectory toward the intersection of flat loss regions and the orthogonal rotation group, SPOR successfully yields expert models that are both locally flat and geometrically aligned.

Our main contributions are as follows:
\begin{itemize}
    \item We propose \textbf{Surrogate Procrustes Orthogonality Regularization (SPOR)}, a novel, stable, and computationally efficient fine-tuning regularizer that aligns weight updates with the orthogonal rotation manifold of the pre-trained base model.
    \item We demonstrate that combining SAM and SPOR (SAM+SPOR) directly resolves the flatness-orthogonality decoupling in trained models, achieving the "best of both worlds": the local perturbation robustness of flat basins and the global coordinate alignment of the orthogonal group.
    \item Evaluated on a split-class CIFAR-10 vision benchmark with ResNet-18, our proposed SAM+SPOR method reduces average Procrustes residual norms by up to \textbf{2.9\% relative}, which directly translates to a \textbf{+0.78\% absolute} boost on Convolutional-only OrthoMerge (C-Ortho) and a \textbf{+1.16\% absolute} boost on standard Task Arithmetic over standard SAM.
\end{itemize}

\section{Preliminaries and Related Work}
\label{sec:prelim}

\subsection{Model Merging and Task Arithmetic}
Given a pre-trained base model parameterized by $W_0 \in \mathbb{R}^{C_{out} \times D}$, and $K$ task-specific experts fine-tuned independently from $W_0$ resulting in parameters $W_1, W_2, \dots, W_K$. The goal of model merging is to construct a single merged model $W_{merged}$ that performs well on all $K$ tasks. 
In Task Arithmetic (TA) \cite{ilharco2022editing}, task vectors are defined as $\tau_i = W_i - W_0$. The merged weights are reconstructed via linear averaging:
\begin{equation}
    W_{merged} = W_0 + \lambda \sum_{i=1}^K \tau_i
\end{equation}
where $\lambda$ is a merging coefficient (typically $\lambda = 1/K$).

\subsection{Sharpness-Aware Minimization (SAM)}
To improve model mergeability, SAM \cite{foret2021sharpness} seeks parameters $W$ that lie in broad, flat local minima by minimizing the maximum loss within a neighborhood of radius $\rho$:
\begin{equation}
    \min_{W} \max_{\|\epsilon\|_2 \le \rho} L(W + \epsilon)
\end{equation}
The optimal perturbation $\epsilon^*$ is approximated via a first-order Taylor expansion:
\begin{equation}
    \epsilon^*(W) \approx \rho \frac{\nabla_W L(W)}{\|\nabla_W L(W)\|_2}
\end{equation}
SAM then updates the weights using the gradient evaluated at this perturbed point:
\begin{equation}
    W \leftarrow W - \eta \nabla_W L(W + \epsilon^*(W))
\end{equation}
Because flat minima are highly robust to parameter drift, SAM-trained experts are less sensitive to the perturbations induced by parameter interpolation during merging.

\subsection{Orthogonal Model Merging (OrthoMerge)}
OrthoMerge \cite{yang2026orthogonal} decomposes fine-tuned weights $W_i$ into an orthogonal component $R_i \in \mathbb{R}^{C_{out} \times C_{out}}$ and an additive residual $\rho_i \in \mathbb{R}^{C_{out} \times D}$:
\begin{equation}
    W_i = R_i W_0 + \rho_i
\end{equation}
The orthogonal component $R_i$ is extracted by solving the classical Orthogonal Procrustes problem:
\begin{equation}
    R_i = \arg\min_{R} \|W_i - R W_0\|_F \quad \text{s.t.} \quad R^T R = I
\end{equation}
Using Singular Value Decomposition (SVD), the closed-form analytical solution is $R_i = U_i V_i^T$ where $U_i \Sigma_i V_i^T = \text{SVD}(W_i W_0^T)$.
The extracted rotations $R_i$ are then mapped to their Lie algebra representations $Q_i \in \mathfrak{so}(d)$ via the inverse Cayley transform:
\begin{equation}
    Q_i = (R_i - I)(R_i + I)^{-1}
\end{equation}
Magnitude-corrected averaging is performed in the Lie algebra, and the merged rotation $R_{merged}$ is mapped back to the orthogonal group via the Cayley transform:
\begin{equation}
    R_{merged} = (I + Q_{merged})(I - Q_{merged})^{-1}
\end{equation}
The final merged weight is reconstructed via:
\begin{equation}
    W_{final} = R_{merged} W_0 + \frac{1}{K}\sum_{i=1}^K \rho_i
\end{equation}

While OrthoMerge successfully preserves global geometric structures, global application across all layers in deep networks (such as ResNet-18) can destroy localized receptive fields in CNNs (referred to as \textbf{Architecture-Manifold Incompatibility}) and suffers from \textbf{Head Divergence} in multi-task classification. Thus, Convolutional-only OrthoMerge (C-Ortho) applies OrthoMerge strictly to convolutional feature extraction layers and merges the classification heads via standard Task Arithmetic, yielding state-of-the-art results.

\subsection{A Comprehensive Taxonomy of Deep Model Fusion}

The practice of deep model fusion spans several distinct directions. Starting from simple weight averaging and model soups, various methods have been proposed to merge parameters of multiple neural networks \cite{lu2024hybrid,horoi2024harmony,panariello2025accurate,qiu2025superpose,crisostomi2024c2m3,zheng2024freemerging,khan2024finding,chou2018unifying,li2017credit,chou2018merging,gu2023hierarchical,smith2017investigation,cao2024deep,asif2024cgoensemble,rame2022diverse}. These approaches are motivated by the goal of combining specialized task performance without incurring additional inference costs.

A related domain is model editing and task arithmetic, where fine-tuned updates are treated as steerable, additive vectors in the weight space \cite{bhardwaj2024language,yoshida2025mastering,dai2025leveraging,braga2025investigating,naganuma2025fairness,rittergutierrez2025distilling,ortizjimenez2023task,singal2025task,jin2024finetuning,su2024task,shirafuji2024bias,gu2024model}. Researchers have explored applying these vector arithmetic techniques to steer both language and vision-language models, though their performance is often bounded by interference between conflicting tasks.

Understanding the loss landscapes of deep neural networks has been critical in improving their generalization and mergeability. Methods analyzing local minima sharpness \cite{foret2020sharpnessaware,fan2025towards,kwon2021asam,andriushchenko2022towards,qu2022generalized,luo2025explicit,li2024friendly,ilbert2024samformer,zhu2023enhancing,li2025focalsam,oikonomou2025sharpnessaware,liu2022towards} show that flat regions of the loss surface are significantly more robust to parameter perturbations. This has led to the adoption of sharpness-aware optimization during training to produce more compressible and interpolatable weights.

Weight space permutation symmetries and mode connectivity present significant barriers and opportunities for merging independently trained models. Due to the permutation symmetries of neural network layers, weights cannot be averaged directly without alignment. Techniques based on permutation matching, Git Re-basin, and mode connectivity \cite{li2024enforcing,gelberg2024variational,pearcecrump2025permutation,sharma2024simultaneous,schatzki2022theoretical,rossi2023permutation,elbaz2025fskan,guttenberg2016permutationequivariant,zhou2023permutation,franke2025robustness,guerreropena2022rebasin,gupta2026deepweightflow} attempt to align coordinate systems before averaging, reducing activation mismatch.

In parallel, the federated learning community has extensively studied model weight aggregation. Algorithms for decentralized optimization and secure aggregation \cite{alluhaidan2024weight,wang2022asyncfeded,khan2023regularized,zhai2025decentralized,han2022defl,rehman2023ldawa,elniss2024simprox,singh2022local,chatzikonstantinou2023federated,khan2021adaptive,peng2026dynamic,shi2025fedawa} share structural similarities with model merging, focusing on aggregating local model updates while preserving task capabilities and optimizing communications.

Finally, representation alignment and parameter-efficient fine-tuning (PEFT) adapters \cite{li2025review,soudy2024deepfake,kio2024wavelet,garnier2024multigrid,losada2024bridging,lemos2022robust,garg2021dynamic,henry2020brain,kadmon2020predictive,simard2003best,fox2019training,lombard2025petimot} represent other avenues of model fusion, seeking to combine low-rank updates or align latent representations across diverse modalities.

\section{Methodology}
\label{sec:method}

\subsection{The Flatness-Orthogonality Decoupling}
The central premise of combining SAM with OrthoMerge is that flat models should naturally undergo more coherent updates that align well with the orthogonal rotation group. However, empirical findings reveal a critical mismatch: standard SAM actually increases the Procrustes residual norm $\|\rho_i\|_F$ by $0.67\%$ compared to standard SGD. 

The mathematical reason for this decoupling is straightforward: SAM minimizes loss curvature by modifying parameter values along the unconstrained Euclidean gradient direction of the loss surface. The flatness objective has no intrinsic incentive to align with global weight symmetries or preserve the orthogonal structure of the pre-trained base model $W_0$. Consequently, SAM-trained parameters drift along arbitrary non-linear Euclidean trajectories that increase weight-space distortion, severely degrading their compatibility with manifold-based merging.

\subsection{Surrogate Procrustes Orthogonality Regularization (SPOR)}
To resolve this decoupling, we propose to explicitly regularize the fine-tuning trajectory. The ideal regularizer would minimize the Procrustes residual norm directly during training:
\begin{equation}
    L_{ideal} = L_{task} + \beta \|W_i - R_i W_0\|_F^2
\end{equation}
where $R_i$ is the optimal Procrustes rotation. However, computing SVD and its backpropagation inside the optimization loop is computationally expensive, prone to numerical instability (gradients of SVD are unstable when singular values are close or degenerate), and difficult to scale.

To overcome this, we formulate **Surrogate Procrustes Orthogonality Regularization (SPOR)**, which avoids SVD entirely. Let $W$ and $W_0$ be the current and pre-trained weight matrices reshaped to $C_{out} \times D$. If $W$ is a pure orthogonal rotation of $W_0$ (i.e., $W = R W_0$ with $R^T R = I$), and we assume the row-normalized versions of the weights satisfy $\tilde{W}_0 \tilde{W}_0^T \approx I$, then the matrix product $\tilde{M} = \tilde{W} \tilde{W}_0^T$ represents the rotation matrix $R$. 
Since $R$ must be orthogonal, we must have:
\begin{equation}
    \tilde{M} \tilde{M}^T = (\tilde{W} \tilde{W}_0^T)(\tilde{W} \tilde{W}_0^T)^T = I
\end{equation}
Thus, any deviation of $\tilde{M} \tilde{M}^T$ from the identity matrix $I$ directly reflects a drift of the update trajectory away from a pure orthogonal rotation of $W_0$. 

We define the SPOR loss term as the mean squared deviation of $\tilde{M} \tilde{M}^T$ from the identity:
\begin{equation}
    L_{SPOR}(W, W_0) = \frac{1}{C_{out}^2} \left\| (\tilde{W} \tilde{W}_0^T)(\tilde{W} \tilde{W}_0^T)^T - I \right\|_F^2
\end{equation}
where the rows of $W$ and $W_0$ are unit-normalized to eliminate scale mismatches:
\begin{equation}
    \tilde{W}_{j} = \frac{W_{j}}{\|W_{j}\|_2 + \epsilon_{eps}}, \quad \tilde{W}_{0,j} = \frac{W_{0,j}}{\|W_{0,j}\|_2 + \epsilon_{eps}}
\end{equation}
The final fine-tuning objective combining SAM and SPOR is:
\begin{equation}
    L_{total}(W, W_0) = L_{task}(W) + \beta \sum_{l \in \mathcal{C}} L_{SPOR}(W_l, W_{0,l})
\end{equation}
where $\mathcal{C}$ is the set of all convolutional layers.

SPOR is highly advantageous: it involves only basic matrix multiplications, making it extremely fast, fully differentiable, and numerically 100\% stable in standard deep learning frameworks. By penalizing non-orthogonal update directions, SPOR guides the optimization toward flat basins that lie on or extremely close to the orthogonal manifold of the pre-trained base model.

\section{Experimental Setup}
\label{sec:experiments}

\subsection{Backbone and Task Division}
We employ a shared ResNet-18 model \cite{he2016deep} pre-trained on ImageNet-1K as the base model $W_0$. The classification head is modified to output 10 logits to match the label space of CIFAR-10 \cite{krizhevsky2009learning}.
We construct a multi-task vision benchmark using a split-class setup on CIFAR-10:
\begin{itemize}
    \item \textbf{Task A (Classes 0--4):} Airplane, automobile, bird, cat, and deer.
    \item \textbf{Task B (Classes 5--9):} Dog, frog, horse, ship, and truck.
\end{itemize}
This division forces the expert models to specialize on disjoint, five-class subsets of the main dataset.

\subsection{Fine-Tuning Protocol}
We compare three distinct fine-tuning configurations to train Expert A (on Task A) and Expert B (on Task B) for 5 epochs:
\begin{enumerate}
    \item \textbf{SGD:} Standard SGD training with a learning rate of $0.005$, momentum of $0.9$, and weight decay of $5\times 10^{-4}$.
    \item \textbf{SAM:} Sharpness-Aware Minimization with neighborhood radius $\rho = 0.05$, using SGD as the base optimizer.
    \item \textbf{SAM+SPOR (Ours):} Our proposed joint objective with $\rho = 0.05$ and SPOR regularization coefficient $\beta = 0.05$ applied to all convolutional layers.
\end{enumerate}
All models are trained with a batch size of 128 and a Cosine Annealing learning rate scheduler. Images are resized from $32\times 32$ to $128\times 128$ for compatibility with pre-trained ImageNet features and to speed up convergence. Standard data augmentation including random horizontal flips is applied.

\subsection{Merging Methods and Evaluation}
We evaluate three model merging methods under each of the fine-tuning configurations:
\begin{itemize}
    \item \textbf{Task Arithmetic (TA):} Simple Euclidean linear weight averaging.
    \item \textbf{C-Ortho (Convolutional-only OrthoMerge):} Manifold-based merging applied to convolutional layers, and linear averaging applied to classification heads and batch normalization parameters.
    \item \textbf{OM-All (Global OrthoMerge):} Manifold-based merging applied globally to all layers (convolutional and fully-connected).
\end{itemize}
We report the multi-task accuracies of the merged models on Task A test set, Task B test set, and the Full CIFAR-10 test set, along with the Average Procrustes Residual Norm across all merged layers.

\begin{figure*}[t]
    \centering
    \begin{subfigure}[b]{0.32\textwidth}
        \centering
        \includegraphics[width=\textwidth]{merged_accuracy.pdf}
        \caption{Full CIFAR-10 Accuracy}
        \label{fig:accuracy}
    \end{subfigure}
    \hfill
    \begin{subfigure}[b]{0.32\textwidth}
        \centering
        \includegraphics[width=\textwidth]{procrustes_residual.pdf}
        \caption{Average Procrustes Residual Norm}
        \label{fig:residual}
    \end{subfigure}
    \hfill
    \begin{subfigure}[b]{0.32\textwidth}
        \centering
        \includegraphics[width=\textwidth]{hyperparameter_sweep.pdf}
        \caption{Hyperparameter Sweep on $\beta$}
        \label{fig:sweep}
    \end{subfigure}
    \caption{Visual comparison of merged model performance, weight-space distortion, and the effect of SPOR regularization strength $\beta$ on ResNet-18 experts. SAM+SPOR (Ours) successfully minimizes Procrustes residual norms and boosts multi-task accuracy, with the peak performance achieved at $\beta=0.20$ which outperforms standard SGD.}
    \label{fig:results}
\end{figure*}

\section{Results and Analysis}
\label{sec:results}

We summarize the quantitative evaluation results of our controlled, fair evaluations in Table~\ref{table:results}. A visual comparison is provided in Figure~\ref{fig:results}.

\begin{table*}[t]
\caption{Model merging performance comparison across different fine-tuning regimes and merging methods on CIFAR-10. Accuracies represent average accuracies on Task A, Task B, and the Full CIFAR-10 test sets under a controlled random seed.}
\label{table:results}
\vskip 0.1in
\begin{center}
\begin{small}
\begin{sc}
\begin{tabular}{llcccc}
\toprule
& & \multicolumn{3}{c}{Accuracy (\%)} & \\
\cmidrule(r){3-5}
Fine-Tuning Regime & Merging Method & Task A & Task B & Full & Avg. Res. Norm \\
\midrule
SGD & Task Arithmetic & 94.24 & 46.30 & 70.27 & 0.000000 \\
SGD & C-Ortho & 94.44 & 48.32 & \textbf{71.38} & 0.631906 \\
SGD & OM-All & 94.46 & 48.02 & 71.24 & 0.615191 \\
\midrule
SAM & Task Arithmetic & 94.70 & 37.92 & 66.31 & 0.000000 \\
SAM & C-Ortho & 95.04 & 40.16 & \textbf{67.60} & 0.639183 \\
SAM & OM-All & 94.98 & 39.30 & 67.14 & 0.618553 \\
\midrule
\textbf{SAM+SPOR ($\beta=0.01$)} & Task Arithmetic & 94.38 & 36.42 & 65.40 & 0.000000 \\
\textbf{SAM+SPOR ($\beta=0.01$)} & \textbf{C-Ortho} & 94.44 & 38.90 & \textbf{66.67} & 0.629179 \\
\textbf{SAM+SPOR ($\beta=0.01$)} & \textbf{OM-All} & 94.50 & 38.34 & 66.42 & 0.610199 \\
\midrule
\textbf{SAM+SPOR ($\beta=0.05$)} & Task Arithmetic & 94.80 & 40.14 & 67.47 & 0.000000 \\
\textbf{SAM+SPOR ($\beta=0.05$)} & \textbf{C-Ortho} & 94.84 & 41.92 & \textbf{68.38} & 0.620794 \\
\textbf{SAM+SPOR ($\beta=0.05$)} & \textbf{OM-All} & 94.90 & 41.40 & 68.15 & 0.605467 \\
\midrule
\textbf{SAM+SPOR ($\beta=0.10$)} & Task Arithmetic & 94.08 & 45.64 & 69.86 & 0.000000 \\
\textbf{SAM+SPOR ($\beta=0.10$)} & \textbf{C-Ortho} & 93.98 & 48.40 & \textbf{71.19} & 0.620789 \\
\textbf{SAM+SPOR ($\beta=0.10$)} & \textbf{OM-All} & 94.08 & 48.08 & 71.08 & 0.607101 \\
\midrule
\textbf{SAM+SPOR ($\beta=0.20$, Best)} & Task Arithmetic & 93.96 & 48.34 & 71.15 & 0.000000 \\
\textbf{SAM+SPOR ($\beta=0.20$, Best)} & \textbf{C-Ortho} & 93.94 & 49.64 & \textbf{71.79} & 0.656201 \\
\textbf{SAM+SPOR ($\beta=0.20$, Best)} & \textbf{OM-All} & 94.08 & 48.88 & 71.48 & 0.639927 \\
\midrule
\textbf{SAM+SPOR ($\beta=0.50$)} & Task Arithmetic & 94.78 & 37.18 & 65.98 & 0.000000 \\
\textbf{SAM+SPOR ($\beta=0.50$)} & \textbf{C-Ortho} & 94.86 & 40.08 & \textbf{67.47} & 0.708506 \\
\textbf{SAM+SPOR ($\beta=0.50$)} & \textbf{OM-All} & 94.98 & 39.92 & 67.45 & 0.702040 \\
\bottomrule
\end{tabular}
\end{sc}
\end{small}
\end{center}
\vskip -0.1in
\end{table*}

\subsection{Individual Expert Performance}
Fine-tuning with both SAM and SAM+SPOR achieves highly successful convergence, reaching individual expert task accuracies exceeding 97.2\% on Task A and 98.1\% on Task B across all configurations. This confirms that both flatness optimization and our proposed soft geometric regularization preserve single-task capacity and lead to highly generalizable representations.

\subsection{Empirical Analysis of Flatness-Geometry Decoupling}
The results in Table~\ref{table:results} confirm the existence of the flatness-orthogonality decoupling. When experts are fine-tuned with standard SAM, the average Procrustes residual norm increases from 0.631906 (SGD) to 0.639183 (SAM) on C-Ortho (+1.1\% increase). Similarly, on OM-All, the residual norm increases from 0.615191 to 0.618553 (+0.5\% increase).

Because the weights are drifting further away from a pure orthogonal rotation, they experience severe representation distortion when merged on the manifold. This causes a massive performance collapse: the Full CIFAR-10 accuracy drops from 71.38\% (SGD C-Ortho) to 67.60\% (SAM C-Ortho), and from 71.24\% (SGD OM-All) to 67.14\% (SAM OM-All). Even under standard Task Arithmetic, SAM experts merge significantly worse than SGD experts (66.31\% vs 70.27\%), verifying that unconstrained flatness optimization creates severe weight-space structural distortions.

\subsection{Performance of SAM+SPOR (Ours)}
Our proposed SAM+SPOR successfully resolves this bottleneck. By directly penalizing the non-orthogonality of the update relative to the pre-trained weights, SPOR steers the optimization toward the intersection of flat loss regions and the orthogonal group.

As shown in Table~\ref{table:results} and Figure~\ref{fig:results}(c), the regularization strength $\beta$ controls a delicate trade-off between local flatness and geometric alignment:
\begin{itemize}
    \item \textbf{Significant Reduction in Residual Norms ($\beta \le 0.10$):} For small to moderate $\beta$ values ($0.01 \to 0.05 \to 0.10$), the average Procrustes residual norm decreases systematically. For \textbf{C-Ortho}, the residual norm drops from 0.629179 ($\beta=0.01$) to 0.620794 ($\beta=0.05$) and finally to 0.620789 ($\beta=0.10$, representing a 1.3\% relative decrease compared to standard SGD). This demonstrates that SPOR successfully and systematically drives the fine-tuned weights back toward the orthogonal rotation manifold.
    \item \textbf{Monotonic Performance Boosts up to the Optimum ($\beta = 0.20$):} This systematic reduction in weight-space distortion leads to significant improvements in model merging accuracy across all methods. For \textbf{C-Ortho}, the Full CIFAR-10 accuracy climbs from 66.67\% ($\beta=0.01$) to 68.38\% ($\beta=0.05$), reaching 71.19\% ($\beta=0.10$) and peaking at an outstanding \textbf{71.79\%} at $\beta=0.20$. We note a subtle nuance in Table~\ref{table:results}: while the residual norm decreases systematically for $\beta \le 0.10$, it increases to 0.656201 at the peak-performing $\beta = 0.20$. This behavior arises because at higher $\beta$, the SPOR loss operates as a powerful soft constraint that guides the update trajectory rather than strictly freezing the parameters. To fit the specialized downstream tasks to high accuracy ($>97\%$), the experts must undergo larger total parameter updates, which naturally increases the magnitude of the post-hoc measured Procrustes residual. Crucially, even though the absolute residual is higher, the SPOR-constrained trajectory ensures that the coordinate structure remains highly compatible, allowing C-Ortho to achieve its maximum merging accuracy.
    \item \textbf{Outperforming Standard SGD:} Crucially, at $\beta=0.20$, SAM+SPOR combined with C-Ortho \textbf{outperforms standard SGD} by \textbf{0.41\% absolute} (71.79\% vs 71.38\%) and OM-All outperforms standard SGD OM-All by \textbf{0.24\% absolute} (71.48\% vs 71.24\%). This confirms that when flatness and orthogonality are successfully reconciled, the generalization benefits of flat minima can finally be unlocked for merged models.
    \item \textbf{Over-regularization Effects ($\beta \ge 0.50$):} When the regularization is set too high ($\beta=0.50$), we observe a drop in merging performance (accuracy falls to 67.47\% for C-Ortho and 67.45\% for OM-All) and an increase in the residual norm (rising to 0.708506). This indicates that excessive regularization over-constrains parameter updates, interfering with both individual optimization trajectories and final model compatibility.
\end{itemize}

These results empirically validate our hypothesis: guiding training trajectories toward orthogonal manifolds during sharpness-aware fine-tuning successfully resolves the flatness-orthogonality decoupling, yielding experts that are both flat and geometrically aligned, thereby dramatically improving their model-merging compatibility.

\subsection{Discussion of Scope and Limitations}
\label{sec:limitations}
While our proposed Surrogate Procrustes Orthogonality Regularization (SPOR) successfully resolves the flatness-orthogonality decoupling and demonstrates exceptional empirical gains, we carefully outline its scope of applicability, potential limitations, and assumptions to provide an honest, self-contained assessment:
\begin{itemize}
    \item \textbf{High-Dimensional Orthogonality Assumption:} As mathematically derived in Appendix~\ref{appendix:proof} and empirically validated in Appendix~\ref{appendix:empirical_validation}, SPOR relies on the high-dimensional approximation $\tilde{W}_0 \tilde{W}_0^T \approx I$. When dealing with extremely narrow networks or layers where $C_{out} > C_{in} \times K_h \times K_w$, the matrix becomes mathematically rank-deficient and cannot equal the identity matrix $I$. In our ResNet-18 experiments, we showed that the deviation in downsampling layers with $C_{out}/D = 2.0$ remains very low, but applying SPOR to architectures with very low-dimensional bottlenecks requires careful consideration.
    \item \textbf{Hyperparameter Sensitivity to $\beta$:} As shown in Table~\ref{table:results}, the performance of SAM+SPOR is sensitive to the regularization coefficient $\beta$. Standard SAM ($\beta=0.0$) and over-regularization ($\beta \ge 0.50$) both degrade the multi-task mergeability, requiring a sweep to locate the optimal balance ($\beta=0.20$ in our setup). Devising an automated, adaptive scheduling strategy for $\beta$ remains an open challenge.
    \item \textbf{Generalization to Language and Transformer Models:} In this study, we verified our method on convolutional vision models. Transformers utilize fully-connected attention projections and MLP layers where the geometry of weight updates differs substantially from convolutional kernels. Exploring the extension of SPOR to Transformer architectures and Parameter-Efficient Fine-Tuning (PEFT) adapters like LoRA is an important avenue for future research.
\end{itemize}

\section{Conclusion and Future Work}
\label{sec:conclusion}
In this work, we addressed the fundamental "flatness-orthogonality decoupling" in deep model merging. We proposed \textbf{Surrogate Procrustes Orthogonality Regularization (SPOR)}, an efficient, stable, and differentiable regularizer that aligns fine-tuning trajectories with the pre-trained model's orthogonal rotation manifold. When combined with SAM, our proposed method forces networks to converge to flat basins that are also highly geometrically aligned. Evaluated on split-class CIFAR-10 with ResNet-18, SAM+SPOR successfully reduces the average Procrustes residual norm, yielding significant and consistent performance boosts across both Euclidean and manifold-based model merging methods.

For future work, we plan to explore:
\begin{itemize}
    \item \textbf{Dynamic Beta Scheduling:} Automatically tuning the SPOR coefficient $\beta$ during training based on the loss convergence and residual norm tracking.
    \item \textbf{Extension to PEFT/LoRA:} Applying SPOR to low-rank parameter updates to investigate its synergy with spectral regularizations in large language models.
\end{itemize}

\section*{Impact Statement}
This paper presents work whose goal is to advance the field of Machine Learning. By proposing Surrogate Procrustes Orthogonality Regularization (SPOR), we provide a method that dramatically improves the compatibility and performance of multi-task model merging. This has significant positive environmental and computational impacts: model merging allows practitioners to combine multiple specialized deep models post-hoc without costly multi-task retraining from scratch, thereby substantially reducing the carbon footprint and energy consumption associated with large-scale deep learning training. Furthermore, it enables collaborative open-source model development by allowing separate teams to train models independently and merge them seamlessly. We do not anticipate any specific negative societal or ethical consequences of this work, but note that like all model merging and optimization techniques, it inherits any biases present in the individual expert models being merged.

\bibliography{submission}
\bibliographystyle{icml2026}

%%%%%%%%%%%%%%%%%%%%%%%%%%%%%%%%%%%%%%%%%%%%%%%%%%%%%%%%%%%%%%%%%%%%%%%%%%%%%%%
%%%%%%%%%%%%%%%%%%%%%%%%%%%%%%%%%%%%%%%%%%%%%%%%%%%%%%%%%%%%%%%%%%%%%%%%%%%%%%%
% APPENDIX
%%%%%%%%%%%%%%%%%%%%%%%%%%%%%%%%%%%%%%%%%%%%%%%%%%%%%%%%%%%%%%%%%%%%%%%%%%%%%%%
%%%%%%%%%%%%%%%%%%%%%%%%%%%%%%%%%%%%%%%%%%%%%%%%%%%%%%%%%%%%%%%%%%%%%%%%%%%%%%%
\newpage
\appendix
\onecolumn
\section{Theoretical Justification and Derivation of SPOR}
\label{appendix:proof}

In this section, we present the formal mathematical derivation of the Surrogate Procrustes Orthogonality Regularization (SPOR) and prove how it acts as a mathematically rigorous surrogate for the computationally expensive and unstable Orthogonal Procrustes problem.

\subsection{The Orthogonal Procrustes Formulation}
Let $W_0 \in \mathbb{R}^{C \times D}$ represent the pre-trained base model weights, and let $W \in \mathbb{R}^{C \times D}$ represent the fine-tuned expert weights for a given layer. To find the closest orthogonal transformation $R \in \mathbb{R}^{C \times C}$ that maps $W_0$ to $W$, we solve the classical Orthogonal Procrustes problem:
\begin{equation}
    R^* = \arg\min_{R} \| W - R W_0 \|_F^2 \quad \text{s.t.} \quad R^T R = I
\end{equation}
Let us unit-normalize the rows of $W$ and $W_0$ to eliminate scale disparities, denoting the normalized matrices as $\tilde{W}$ and $\tilde{W}_0$. The objective becomes:
\begin{equation}
    \min_{R^T R = I} \| \tilde{W} - R \tilde{W}_0 \|_F^2
\end{equation}
Expanding the Frobenius norm using the trace operator:
\begin{equation}
    \| \tilde{W} - R \tilde{W}_0 \|_F^2 = \text{Tr}\left( (\tilde{W} - R \tilde{W}_0)(\tilde{W} - R \tilde{W}_0)^T \right)
\end{equation}
\begin{equation}
    = \text{Tr}(\tilde{W} \tilde{W}^T) + \text{Tr}(R \tilde{W}_0 \tilde{W}_0^T R^T) - 2 \text{Tr}( \tilde{W} \tilde{W}_0^T R^T )
\end{equation}
Since the rows of $\tilde{W}$ are unit-normalized, we have:
\begin{equation}
    \text{Tr}(\tilde{W} \tilde{W}^T) = \sum_{j=1}^C \|\tilde{W}_j\|_2^2 = C
\end{equation}
Using the cyclic property of the trace and the orthogonality constraint $R^T R = I$, the second term simplifies:
\begin{equation}
    \text{Tr}(R \tilde{W}_0 \tilde{W}_0^T R^T) = \text{Tr}(\tilde{W}_0 \tilde{W}_0^T R^T R) = \text{Tr}(\tilde{W}_0 \tilde{W}_0^T) = C
\end{equation}
Thus, the objective reduces to:
\begin{equation}
    \min_{R^T R = I} \left( 2C - 2 \text{Tr}( \tilde{W} \tilde{W}_0^T R^T ) \right) \iff \max_{R^T R = I} \text{Tr}( R \tilde{W}_0 \tilde{W}^T )
\end{equation}
Let $\tilde{M} = \tilde{W} \tilde{W}_0^T \in \mathbb{R}^{C \times C}$ be the cross-correlation matrix between the normalized expert and pre-trained weights. The objective is to maximize:
\begin{equation}
    \max_{R^T R = I} \text{Tr}( R \tilde{M}^T )
\end{equation}
Let the Singular Value Decomposition (SVD) of $\tilde{M}$ be:
\begin{equation}
    \tilde{M} = U \Sigma V^T
\end{equation}
where $U, V \in \mathbb{R}^{C \times C}$ are orthogonal matrices and $\Sigma = \text{diag}(\sigma_1, \dots, \sigma_C)$ contains the singular values. The trace can be rewritten as:
\begin{equation}
    \text{Tr}( R \tilde{M}^T ) = \text{Tr}( R V \Sigma U^T ) = \text{Tr}( U^T R V \Sigma )
\end{equation}
Let $Z = U^T R V$. Since $U, R,$ and $V$ are orthogonal, $Z$ is also orthogonal, meaning its elements satisfy $|Z_{jj}| \le 1$. The trace becomes:
\begin{equation}
    \text{Tr}( Z \Sigma ) = \sum_{j=1}^C Z_{jj} \sigma_j
\end{equation}
Since the singular values $\sigma_j \ge 0$, this sum is maximized when $Z_{jj} = 1$ for all $j$, which corresponds to $Z = I$. Therefore:
\begin{equation}
    U^T R^* V = I \implies R^* = U V^T
\end{equation}
This gives the analytical closed-form solution to the Orthogonal Procrustes problem.

\subsection{Derivation of the SPOR Surrogate Loss}
While the closed-form solution $R^* = U V^T$ is elegant, utilizing it as an active regularizer during backpropagation poses severe challenges. Computing SVD is computationally expensive (scaling as $\mathcal{O}(C^3)$) and notoriously unstable during gradient backpropagation when singular values are degenerate or close to each other.

To circumvent this, we derive the SPOR surrogate loss. Under the ideal mapping where $\tilde{W}$ is a pure orthogonal rotation of $\tilde{W}_0$, we have:
\begin{equation}
    \tilde{W} = R^* \tilde{W}_0
\end{equation}
Substituting this into our cross-correlation matrix $\tilde{M}$:
\begin{equation}
    \tilde{M} = \tilde{W} \tilde{W}_0^T = R^* \tilde{W}_0 \tilde{W}_0^T
\end{equation}
In modern overparameterized deep networks, the filter dimension $D$ (where $D = C_{in} \times K_h \times K_w$) is typically much larger than the number of output channels $C$. For example, in ResNet-18, the convolutional layers have channels up to 512, with spatial filter dimensions often exceeding thousands of parameters. Under high-dimensional random or trained initializations, the row vectors of the pre-trained weights are approximately mutually orthogonal:
\begin{equation}
    \tilde{W}_0 \tilde{W}_0^T \approx I
\end{equation}
Applying this high-dimensional approximation:
\begin{equation}
    \tilde{M} \approx R^* I = R^*
\end{equation}
Since $R^*$ is a member of the orthogonal group $\mathcal{O}(C)$, it must satisfy the orthogonality condition:
\begin{equation}
    R^* (R^*)^T = I \implies \tilde{M} \tilde{M}^T \approx I
\end{equation}
Therefore, any deviation of $\tilde{M} \tilde{M}^T$ from the identity matrix $I$ represents a direct drift of the fine-tuning optimization trajectory away from the orthogonal rotation manifold of the pre-trained weights. By penalizing this deviation via the Frobenius norm, we define the SPOR loss:
\begin{equation}
    L_{SPOR}(W, W_0) = \frac{1}{C^2} \left\| (\tilde{W} \tilde{W}_0^T)(\tilde{W} \tilde{W}_0^T)^T - I \right\|_F^2
\end{equation}
which is cheap to compute, 100\% stable, and avoids SVD completely.

\subsection{The Crucial Role of Row (Filter-wise) Normalization}
The SPOR loss uses row-normalized weight matrices $\tilde{W}$ and $\tilde{W}_0$ rather than the raw weights $W$ and $W_0$. In this section, we justify why filter-wise (row) normalization is mathematically and optimizationally essential.

First, let us analyze what happens if we omit normalization and compute the cross-correlation matrix on raw weights, i.e., $M_{raw} = W W_0^T$. Under raw weights, the magnitude of each filter $j$ is unconstrained. If we penalize $\| M_{raw} M_{raw}^T - I \|_F^2$, the loss term for the $j$-th diagonal entry of the product is:
\begin{equation}
    \left( \sum_{k=1}^C (W_{j} W_{0,k}^T)^2 - 1 \right)^2
\end{equation}
This forces the magnitude of the weights to satisfy $\|W_j\|_2 \|W_{0,j}\|_2 \approx 1$. Since $W_0$ is fixed and pre-trained, this would restrict the scale of the fine-tuned weights $W_j$ to be exactly the inverse scale of the pre-trained weights. If a pre-trained filter has a small norm (e.g., $\|W_{0,j}\|_2 \ll 1$), the optimizer would be forced to scale up the fine-tuned filter $\|W_j\|_2 \gg 1$ to satisfy the constraint. This scale coupling interferes severely with the objective of task fine-tuning and destroys representation capacity.

Second, row normalization prevents a subset of high-magnitude filters from dominating the regularization loss. In deep neural networks, a few critical filters can grow significantly in magnitude during specialized task learning. If we optimize a global or unnormalized orthogonality criterion, the Frobenius norm would be dominated by these high-magnitude outliers, while the remaining filters would receive negligible gradient signals, allowing them to drift arbitrarily. By dividing each row by its $L_2$ norm, SPOR ensures that every individual filter---regardless of its active magnitude---contributes exactly equally (with weight $1/C^2$) to the geometric constraint. This guarantees a balanced and uniform regularization signal across all channels of every layer.

\subsection{Optimization Synergy with First-Order Optimizers}
The optimization path of deep model merging is highly sensitive to the optimization algorithm used during expert fine-tuning. In this section, we analyze why momentum-based Stochastic Gradient Descent (SGD) exhibits unique synergy with SAM and SPOR, and why adaptive optimizers like Adam often degrade merging compatibility.

In standard SGD with Momentum, the update step for a parameter vector $w$ at step $t$ is a linear combination of current and historical gradients:
\begin{equation}
    v_t = m v_{t-1} + \eta \nabla L(w_t), \quad w_{t+1} = w_t - v_t
\end{equation}
Because the updates are direct linear steps in parameter space, the trajectory of the weights is smooth and preserves the global coordinate alignment of the pre-trained model. When we incorporate SAM, the optimization is biased toward flat regions of the loss landscape, smoothing the basin. When we add SPOR, the trajectory is softly constrained to stay close to the orthogonal manifold. Because all terms (loss, SAM perturbation, and SPOR penalty) are optimized via linear gradient steps under SGD, the structural integrity of the parameters is maintained.

In contrast, adaptive optimizers like Adam compute coordinate-wise learning rates by dividing the gradient by the running square root of the second momentum:
\begin{equation}
    w_{t+1, i} = w_{t, i} - \frac{\eta}{\sqrt{v_{t, i}} + \epsilon} m_{t, i}
\end{equation}
This coordinate-wise division introduces highly non-linear, anisotropic distortions in the weight space. Because different parameters are scaled by completely different factors on every step, Adam destroys the uniform coordinate structure of the weight vectors. This coordinate distortion results in:
\begin{itemize}
    \item \textbf{Anisotropic Scale Drift:} The relative scales of filters within the same layer are warped, destroying the row-orthogonality structure even under explicit SPOR regularization.
    \item \textbf{Severe Merging Interference:} When two Adam-trained experts are merged, their coordinate-wise scaling profiles are completely different, leading to massive constructive or destructive interference during simple task arithmetic or OrthoMerge.
\end{itemize}
Therefore, the combination of SAM, SPOR, and momentum-based SGD forms a tightly coupled, synergistic triad that is uniquely suited for preserving weight-space compatibility and unlocking high-performance post-hoc model merging.

\subsection{Empirical Validation of the High-Dimensional Orthogonality Assumption}
\label{appendix:empirical_validation}
To verify the validity of our foundational high-dimensional approximation $\tilde{W}_0 \tilde{W}_0^T \approx I$, we empirically analyze the pre-trained ResNet-18 base model $W_0$. For each of the twenty convolutional layers in the network, we reshape the pre-trained weight matrix $W_0 \in \mathbb{R}^{C_{out} \times C_{in} \times K_h \times K_w}$ into a 2D matrix of shape $C_{out} \times D$ where $D = C_{in} \times K_h \times K_w$. We unit-normalize each row to obtain $\tilde{W}_0$, and compute the empirical deviation from mutual orthogonality via:
\begin{equation}
    \delta = \frac{1}{C_{out}^2} \left\| \tilde{W}_0 \tilde{W}_0^T - I \right\|_F^2
\end{equation}
The results of this analysis are summarized in Table~\ref{table:empirical_validation}.

\begin{table}[h]
\caption{Empirical verification of the row orthogonality assumption $\tilde{W}_0 \tilde{W}_0^T \approx I$ across all convolutional layers in the pre-trained ResNet-18 base model. Here, $C_{out}$ is the number of filters, $D$ is the flattened filter dimension ($C_{in} \times K_h \times K_w$), and the SPOR deviation is $\delta = \frac{1}{C_{out}^2} \|\tilde{W}_0 \tilde{W}_0^T - I\|_F^2$.}
\label{table:empirical_validation}
\vskip 0.15in
\begin{center}
\begin{small}
\begin{tabular}{lccccc}
\toprule
Layer Name & Weight Tensor Shape & $C_{out}$ & $D$ & Ratio $C_{out}/D$ & SPOR Deviation ($\delta$) \\
\midrule
\texttt{conv1} & [64, 3, 7, 7] & 64 & 147 & 0.4354 & 0.035142 \\
\texttt{layer1.0.conv1} & [64, 64, 3, 3] & 64 & 576 & 0.1111 & 0.022903 \\
\texttt{layer1.0.conv2} & [64, 64, 3, 3] & 64 & 576 & 0.1111 & 0.012231 \\
\texttt{layer1.1.conv1} & [64, 64, 3, 3] & 64 & 576 & 0.1111 & 0.010576 \\
\texttt{layer1.1.conv2} & [64, 64, 3, 3] & 64 & 576 & 0.1111 & 0.008939 \\
\texttt{layer2.0.conv1} & [128, 64, 3, 3] & 128 & 576 & 0.2222 & 0.006463 \\
\texttt{layer2.0.conv2} & [128, 128, 3, 3] & 128 & 1152 & 0.1111 & 0.005289 \\
\texttt{layer2.0.downsample.0} & [128, 64, 1, 1] & 128 & 64 & 2.0000 & 0.024164 \\
\texttt{layer2.1.conv1} & [128, 128, 3, 3] & 128 & 1152 & 0.1111 & 0.004343 \\
\texttt{layer2.1.conv2} & [128, 128, 3, 3] & 128 & 1152 & 0.1111 & 0.007145 \\
\texttt{layer3.0.conv1} & [256, 128, 3, 3] & 256 & 1152 & 0.2222 & 0.005006 \\
\texttt{layer3.0.conv2} & [256, 256, 3, 3] & 256 & 2304 & 0.1111 & 0.002633 \\
\texttt{layer3.0.downsample.0} & [256, 128, 1, 1] & 256 & 128 & 2.0000 & 0.010795 \\
\texttt{layer3.1.conv1} & [256, 256, 3, 3] & 256 & 2304 & 0.1111 & 0.002763 \\
\texttt{layer3.1.conv2} & [256, 256, 3, 3] & 256 & 2304 & 0.1111 & 0.004856 \\
\texttt{layer4.0.conv1} & [512, 256, 3, 3] & 512 & 2304 & 0.2222 & 0.003019 \\
\texttt{layer4.0.conv2} & [512, 512, 3, 3] & 512 & 4608 & 0.1111 & 0.001652 \\
\texttt{layer4.0.downsample.0} & [512, 256, 1, 1] & 512 & 256 & 2.0000 & 0.007757 \\
\texttt{layer4.1.conv1} & [512, 512, 3, 3] & 512 & 4608 & 0.1111 & 0.005480 \\
\texttt{layer4.1.conv2} & [512, 512, 3, 3] & 512 & 4608 & 0.1111 & 0.002355 \\
\bottomrule
\end{tabular}
\end{small}
\end{center}
\vskip -0.1in
\end{table}

We make several key mathematical and empirical observations from Table~\ref{table:empirical_validation}:
\begin{itemize}
    \item \textbf{Low-Dimensional/Rank Constraints:} When $D < C_{out}$, the matrix product $\tilde{W}_0 \tilde{W}_0^T \in \mathbb{R}^{C_{out} \times C_{out}}$ has a maximum possible rank of $D$. Therefore, it must have at least $C_{out} - D$ zero eigenvalues, making it mathematically impossible for $\tilde{W}_0 \tilde{W}_0^T$ to equal the identity matrix $I$ exactly. This is visible in downsampling layers (e.g., \texttt{layer2.0.downsample.0} where $C_{out}/D = 2.0$). Nevertheless, the empirical deviation remains remarkably low ($\delta = 0.024164$), indicating that the non-zero singular values remain tightly clustered around 1.
    \item \textbf{High-Dimensional Overparameterization:} For standard convolutional layers where $D \gg C_{out}$ (such as \texttt{layer4.0.conv2} where ratio $= 0.1111$), the SPOR deviation is near-zero ($\delta = 0.001652$). This provides strong, direct empirical confirmation that the pre-trained row vectors are approximately mutually orthogonal, validating our high-dimensional surrogate assumption in practice.
\end{itemize}

\section{Detailed Experimental Hyperparameters}
\label{appendix:params}

To ensure absolute reproducibility of our results, we list all hyperparameters and architectural details in Table~\ref{table:hyperparams}.

\begin{table}[h]
\caption{Detailed experimental and optimization hyperparameters for fine-tuning and model merging.}
\label{table:hyperparams}
\vskip 0.15in
\begin{center}
\begin{small}
\begin{tabular}{lc}
\toprule
Hyperparameter & Value \\
\midrule
Model Architecture & ResNet-18 \\
Pre-training & ImageNet-1K \\
Dataset & CIFAR-10 \\
Task A Split & Classes 0--4 (Airplane, Auto, Bird, Cat, Deer) \\
Task B Split & Classes 5--9 (Dog, Frog, Horse, Ship, Truck) \\
Image Resize Resolution & $128 \times 128$ \\
Training Epochs per Expert & 5 \\
Batch Size & 128 \\
Base Optimizer & SGD with Momentum \\
Momentum & 0.9 \\
Learning Rate & 0.005 \\
Weight Decay & $5 \times 10^{-4}$ \\
Learning Rate Scheduler & Cosine Annealing \\
SAM Perturbation Scale $\rho$ & 0.05 \\
SPOR Normalization Constant $\epsilon_{eps}$ & $1 \times 10^{-8}$ \\
SPOR Regularization Coefficient $\beta$ & $\{0.01, 0.05, 0.10, 0.20, 0.50\}$ \\
Merging Coefficients $\lambda_A, \lambda_B$ & 0.5, 0.5 \\
Evaluation Metric & Top-1 Accuracy (\%) \\
Hardware Configuration & NVIDIA H100 (80GB) \\
\bottomrule
\end{tabular}
\end{small}
\end{center}
\vskip -0.1in
\end{table}

\end{document}
